\section{Задача}
\label{sec:task}

Вход --- SMILES молекулы. Оцениваемых треков два: регрессия и классификация. Третий, предсказание позы
лиганда в кармане фермента, откроется в середине соревнования и здесь не разбирается.

В регрессионном треке нужно выдать матрицу $\hat{Y} \in \mathbb{R}^{750\times 4}$: $\pic$ прямого
ингибирования для CYP1A2, CYP2C9, CYP2D6 и CYP3A4. В классификационном --- $\hat{Y} \in \{0,1\}^{750\times 2}$:
метку time-dependent ингибирования для 3A4 и 2D6. Тестовые 750 соединений выданы без меток, отправка
закрывается 3 ноября.

\subsection{Метрика регрессии и что из неё следует}

Метрика --- не средняя абсолютная ошибка. Организаторы публикуют вместе с каждым $\pic$ границы
доверительного интервала байесовской подгонки кривой, и ошибка считается до ближайшей границы, а не до
точечной оценки:
\begin{equation}
\text{ST-RAE} \;=\;
\frac{\sum_i \bigl[(\hat y_i - h_i)_+ + (l_i - \hat y_i)_+\bigr]}
     {\sum_i \bigl[(\bar y - h_i)_+ + (l_i - \bar y)_+\bigr]},
\qquad \bar y = \frac{1}{n}\sum_i y_i .
\label{eq:strae}
\end{equation}

Здесь $l_i$ и $h_i$ --- нижняя и верхняя границы интервала $i$-го соединения, а $(\cdot)_+ = \max(\cdot, 0)$.
Числитель: предсказание, попавшее внутрь полосы, не штрафуется вовсе, а вышедшее наружу штрафуется только
на расстояние до ближайшего края. Знаменатель: та же величина для константного предсказателя, выдающего
среднее по оценочному набору. Отношение единица означает, что модель не лучше предсказания одного числа для всех соединений подряд;
меньше единицы --- лучше.

Из этого следует три вещи, и каждая влияет на архитектуру.

Первое. Задача интервально-цензурированная, и обучаться надо на той же клипованной L1, а не на квадратичной
ошибке. Квадратичная будет тянуть предсказание к центру полосы там, где метрике всё равно, где именно
внутри полосы оно стоит.

Второе. Это отношение сумм, а не среднее отношений. Соединения с широкой полосой почти невесомы в обеих
частях: они дают мало и в числитель, и в знаменатель. Метрику определяют соединения с узкой полосой,
то есть активные, у которых кривая хорошо определена.

Третье. Оптимальная точечная оценка под \eqref{eq:strae} не обязана совпадать с условным средним
предсказательного распределения. Это разбирается в разделе~\ref{sec:decision}.

Знаменатель считается по среднему \emph{того набора}, на котором метрика вычисляется. У нас это среднее
по обучающей подвыборке фермента, на лидерборде --- по 750 тестовым соединениям. Значит наши числа и числа
лидерборда напрямую несопоставимы, даже если модель одна и та же; сопоставимы только разности между
вариантами, посчитанные на одном наборе.

\subsection{Метрика классификации}

Для TDI считается коэффициент корреляции Мэтьюса:
\begin{equation}
\text{MCC} = \frac{TP\cdot TN - FP\cdot FN}{\sqrt{(TP+FP)(TP+FN)(TN+FP)(TN+FN)}} .
\label{eq:mcc}
\end{equation}

Это коэффициент корреляции Пирсона между двумя бинарными векторами, предсказанным и истинным. При
несбалансированных классах он ведёт себя куда честнее точности: предсказать всем «нет» даёт точность
78 процентов и MCC ровно ноль.

MCC не раскладывается по объектам. Вклад одного соединения зависит от того, что
предсказано остальным, потому что все четыре счётчика стоят под общим корнем. Отсюда следует, что порог
0.5 почти никогда не оптимален и его надо выбирать явно --- см. раздел~\ref{sec:decision}.

Обе метрики агрегируются макро-усреднением по эндпойнтам, и лидерборд сортируется по этому среднему.
Организаторы пересчитывают его на 1000 бутстрап-выборках.
